%! Graphing Polynomial Functions and Modeling — Unit 4, Lesson 4.6 — Warm-Up (Answer Key)
\documentclass[10pt]{article}
\usepackage{algebra2-article}
\usepackage{algebra2-key}   % pulls in -boxes; do NOT also load -boxes
\newcommand{\TallMath}[1]{$\displaystyle #1\rule[-1.4em]{0pt}{3.2em}$}

\begin{document}

\pageheader{Unit 4, Lesson 4.6}{Warm-Up --- Answer Key}
\namedateperiod

\vspace{0.08in}

\begin{notesbox}{Getting Ready: Three Pieces of One Picture}
\small You already own every piece of a polynomial's graph. Today they get assembled. Each item
below is a tool you have used before --- notice what part of the picture it controls.

\vspace{5pt}
\textbf{1. The arms} (Lessons 4.0 and 4.1). You do \emph{not} have to multiply these out: degrees
add and leading coefficients multiply.

\vspace{3pt}
\renewcommand{\arraystretch}{1.5}
\begin{tabularx}{\linewidth}{|X|C{1.6cm}|C{2.0cm}|C{2.6cm}|C{2.6cm}|}
\hline
\rowcolor{forestbg}
\textbf{Function} & \textbf{Degree} & \textbf{Leading coeff.} & \textbf{As $x \to -\infty$, $f(x) \to$} & \textbf{As $x \to +\infty$, $f(x) \to$} \\ \hline
(a) $f(x) = (x-1)^2(x+2)$   & \ans{$3$} & \ans{$1$}  & \ans{$-\infty$} & \ans{$+\infty$} \\ \hline
(b) $g(x) = -2x(x+3)^3$     & \ans{$4$} & \ans{$-2$} & \ans{$-\infty$} & \ans{$-\infty$} \\ \hline
(c) $h(x) = -x^2(x-3)^3$    & \ans{$5$} & \ans{$-1$} & \ans{$+\infty$} & \ans{$-\infty$} \\ \hline
\end{tabularx}

\vspace{6pt}
\textbf{2. The $x$-intercepts, and what happens at each one} (Lesson 4.0). Here is the polynomial
you have met in every lesson of this unit: \; $g(x) = x^3 - 3x + 2 = (x-1)^2(x+2)$.

\vspace{3pt}
Zeros: $x = $ \ans{$1$} with multiplicity \ans{$2$}, \quad and $x = $ \ans{$-2$} with
multiplicity \ans{$1$}.

\vspace{3pt}
At a zero of \emph{even} multiplicity the graph \ans{touches} the $x$-axis; at a zero of
\emph{odd} multiplicity it \ans{crosses} it.

\vspace{3pt}
The $y$-intercept: $g(0) = $ \ans{$2$}, so the graph passes through the point \ans{$(0,2)$}.

\vspace{6pt}
\textbf{3. Above or below?} The two zeros cut the $x$-axis into three stretches. Test one $x$-value
from each and record only the \emph{sign} of the answer.

\vspace{3pt}
\renewcommand{\arraystretch}{1.5}
\begin{tabularx}{\linewidth}{|C{3.2cm}|C{2.4cm}|C{2.6cm}|X|}
\hline
\rowcolor{forestbg}
\textbf{Interval} & \textbf{Test $x$} & \textbf{$g(x) = $} & \textbf{Sign of $g(x)$ --- is the graph above or below the $x$-axis?} \\ \hline
$x < -2$        & $-3$ & \ans{$-16$} & \ans{negative --- below} \\ \hline
$-2 < x < 1$    & $0$  & \ans{$2$}   & \ans{positive --- above} \\ \hline
$x > 1$         & $2$  & \ans{$4$}   & \ans{positive --- above} \\ \hline
\end{tabularx}

\vspace{4pt}
\textbf{The hinge.} You now know which way both arms point, where the graph meets the $x$-axis and
whether it crosses or touches there, and whether it sits above or below the axis on every stretch in
between. \textbf{Could you pick this graph out of a lineup?} \ans{yes}

\vspace{2pt}
What is the one thing this information still does \emph{not} tell you? \ans{how \emph{high} or
\emph{low} the graph turns --- the exact peaks and valleys}

\end{notesbox}

\begin{teachernote}
Every item here is a tool students already have; the point of the page is that each one controls a
different part of one picture. Item 1 is the 4.0 end-behavior rule plus the 4.1 shortcut (degrees
add, leading coefficients multiply) --- watch for students who expand $(x+3)^3$ by hand, and for
$-2x(x+3)^3$ read as degree $3$. Item 2 is multiplicity from 4.0. Item 3 is new only in packaging:
testing one point per interval is what a \textbf{sign chart} does, and item 3 \emph{is} the sign
chart for the unit anchor.

\vspace{3pt}
\textbf{Debrief by assembling, not by checking.} Put the three results on the board together ---
arms down-left / up-right, a touch at $1$ and a crossing at $-2$, below then above then above ---
and ask whether anything about the picture is still undecided. Accept ``how high the hump is.'' That
is the honest answer: the five-step plan in the notes locates \emph{every} feature except the exact
coordinates of the turning points, which need technology or calculus. Say so out loud; it keeps
students from thinking they did something wrong.

\vspace{3pt}
Common wrong turns: end behavior taken from the \emph{constant} instead of the leading term; the
$(0,2)$ point written as $(2,0)$; and sign work done by re-expanding $x^3-3x+2$ instead of using the
factored form, where $g(-3) = (-4)^2(-1)$ is a ten-second calculation. Push the factored form --- it
is the whole reason the sign chart is fast.
\end{teachernote}

\end{document}
